\cotefile{53}
\scanpage{9}
% slanted
\margin{C'est \unsure{faux|sans|sous}} \deleted{\unsure{Coroll.|Coroll 1}} \uline{Corollaire} 2 \deleted{\unsure{Toute|Toutes}} \inserted{Les} \uline{composantes irréductibles de $G$} \deleted{(\unsure{supposé} loc. de t. f.)} \uline{sont les composantes connexes}. \uline{En particulier}, $G^{\unsure{0|\circ}}$ \uline{est irréductible}.

On peut supposer $k$ parfait, $G$ réduit donc normal (resp. lisse). Mais alors c'est \uline{vrai pour tt schéma normal} \inserted{(lisse)} sur $k$.... On a supposé ici $G$ loc. de t. f. sur $k$. Sinon on peut dire que $G_0$, la comp. irréd. de \unsure{$e$|$k$}, est stable par générisation, donc contient $\operatorname{\unsure{Loc|loc|Irr}}_e(G)$.

% slanted
\margin{\textcircled{?} \deleted{\unsure{C'est dans|C'est donc}} le cas loc. de t. f. au \unsure{moins}}

% slanted
\margin{$\unsure{-|\leftarrow}$ \unsure{$W^{\mathrm{no}}$|$W^{\mathrm{ro}}$|$N^{\mathrm{ro}}$} en $G$ \uline{loc. de t. f.}}

\deleted{\uline{Corollaire} 3 Pour \deleted{deux} ouverts $U$\inserted{, $V$} de $G$, \inserted{\deleted{\unsure{loc. de t. f.|conn. de t. f.}} On peut \unsure{s'en} \inserted{dispenser} \deleted{supposer} $U$ \deleted{et} $V$ \unsure{qu. cpct|qu. qpct} \uline{ou} $G$ loc. de t. f.)} $U\deleted{\unsure{U'|U}}V$ \deleted{contient $G^0$, si $U,U'$ sont des ouverts de $G$ qui contiennent $e$, $UV\supset G^0$.} \inserted{est une partie à la fois ouverte et fermée de $G$ (\unsure{évidemt|évidemment|évident} qu. cpct)}, donc stable par translation par $G^0$.}

\deleted{\deleted{\unsure{La deuxième assertion|La deuxième étant} conséquence immédiate de 2} \unsure{première|premier|pour voir|pour raison} \deleted{\unsure{on le|on la} \unsure{pour|peut}} \unsure{l'appliquant|l'expliquant} : $W=U\cap U^{-1}\cap V\cap V^{-1}$. \deleted{\illegible{2}} Pour prouver $UU^{-1}\supset G^0$, \deleted{\unsure{on peut}} \inserted{OPS \deleted{\unsure{par}} descente que $U$ contient pt rat $/k$, et par translation que $U\ni e$.} Soit $g\in G^0$, il faut montrer que (sur une ext $k'$ de \unsure{$k(g)$|$\kappa(g)$}) $g=\deleted{\unsure{g|g'}}uv^{-1}$, avec $u,v\in U(k')$, i.e. $(gU\cap U)(k')\ne \unsure{\varnothing|\phi}$, \deleted{\unsure{et} il suffit de prouver que (sur $k_1=k(g)$)} \deleted{$\unsure{gU\cap U\ne\varnothing|gU\cap U\ne e}$} or $gU\cap G^0$ et $U\cap G^0$ sont des ouverts non vides de $G^0$ (contient $g$ et $e$, donc ont une intersection non vide, puisque $G^0$ est irréd.}

\deleted{\uline{Autre dém} $G\times G \unsure{-|\longrightarrow} G$ est plat (car isom à $G\times G\xrightarrow{\mathrm{pr}_{\unsure{1|2}}}G$) \unsure{donc|de} \deleted{\unsure{ses|sa}} \unsure{générisant|générisants}, donc l'image de $U\times \unsure{U|V}$ est stable par générisation, donc si $G$ est loc. de t. f.}

\deleted{La $=$ démonstration montre que $U.V$ \unsure{stable} par translation par $G^0$, donc étant qu. cpct il est \deleted{stable} ouvert et fermé.}

\deleted{$U.V$ ouvert comme image de $U\times_k V\subset G\times_k G$ par $G\times_k G \unsure{-|\longrightarrow} G$, qui est \unsure{ouvert|ouvert(e)} (isom : $G\times G\underset{\mathrm{pr}_1}{\unsure{-|\longrightarrow}}G$ qui l'est).}

\deleted{D'autre part, pour prouver que $U.V$ est fermé, il suffit, comme $U.V$ est \deleted{\unsure{vide ou quasi|vide ou qu. cpct}} qu. cpt donc \unsure{rétrocompact|rétro-compact} (dans $G$ \uline{séparé}), \deleted{pour} de voir que $x\in U.V$, \deleted{$\unsure{\Rightarrow x\in U.V.|\Rightarrow y\in U.V.}$} \inserted{$y\in \unsure{\overline{(x)}|(x)|\overline{\{x\}}}$, alors $y\in U.V$.} On est ramené à \deleted{prouver} le prouver pour $x$}

% archive sheet number: 8
\scanpage{10}
Rectificatif sur critère d'irréductibilité géom.

\diagram{\begin{tikzcd}[ampersand replacement=\&,column sep=1.6em,row sep=.45em]
{} \&[-.3em] {} \&[-.5em] Z'\arrow[dr,dash,"{f'}"'] \& {}\\
Z\arrow[r,"f"] \& X\arrow[d,dash] \& {} \& X'\arrow[ll,"p"']\arrow[d,dash]\\
{} \& k\arrow[rr,dash] \& {} \& k'\\
{} \& {} \& {} \& \deleted{\text{\unsure{sép clos|séparable|sépar. clos}}}
\end{tikzcd}}
Supposons a) $Z$ géom. irréd, b) $f(Z)$ non contenu dans l'ens des points \deleted{\illegible{1}} où passent plusieurs comp. irréd. géom (p. ex. $f(Z)$ contient un pt géom normal, \deleted{\unsure{ou|un}} \inserted{voire} géom régulier, de $X_{\mathrm{red}}$). Alors $X$ géom irréductible.

\uline{Dém} [OPS $Z$ réduit à un pt \deleted{\unsure{de}} $x$ de $X$.] \deleted{\unsure{Alors $X$ irréd|On a $X$ irréd|Alors $X$ étant}} Soit $k'$ \deleted{clôture séparable} \inserted{ext finie \unsure{gal.|galois.}} de $k$, \deleted{\unsure{Alors|Soit}} $X'_i$ une composante irréductible de $X'$.  \deleted{\unsure{Elle domine|Elle devient}} contenant $f'(Z')$\inserted{(qui est \unsure{irréd|irrég})}, les \unsure{transformés|transformées} par $\pi$ sont des composantes irréductibles contenant $f'(Z')$, donc sont \uline{égales} à $\deleted{\unsure{f(|g(}}X'_i$.  Donc $X'_i$ est l'image inverse de une partie fermée irréductible de $X$, savoir de $\unsure{p|p'|f}(X'_i)$, or $p(X'_i)=X$ [platitude \deleted{\unsure{de|du}} \unsure{]|)} $p$ + $p$ fermé], donc $X'_i=X'$, donc $X'$ est irréd.

\unsure{Contre'exemple|Contr'exemple} si on laisse tomber b)

\[
X=\mathbb{R}[x,y]/(x^{\unsure{2|7}}+y^2),\quad Z=\operatorname{Spec}\mathbb{R},
\]
\diagram{\begin{tikzcd}[ampersand replacement=\&,row sep=.2em,column sep=1.5em]
Z\arrow[r] \& X\\
\mathbb{R} \& \mathbb{R}[x,y]/{-}\arrow[l]\\
0 \& x,y\arrow[l,maps to]
\end{tikzcd}}
\par\noindent\hfill\rule{.25\linewidth}{.4pt}\hfill\strut\par
\deleted{\unsure{Complément|Compléments} sur $\pi_0(X/k)$}

\uline{Remarque} Dans l'équivalence $\mathrm{Et}(k)\simeq(\pi\text{-}\mathrm{Ens})$, \deleted{les \unsure{ensembles}} \deleted{\unsure{on voit|on a vu}} \margin{\deleted{$\unsure{k'\underset{\pi}{\longleftarrow}k|k'\underset{\pi}{-}k'}$}}

\diagram{\begingroup\small\(\left\{\begin{array}{@{}l@{\;\;}l@{}}
\multicolumn{1}{c}{X} & \multicolumn{1}{c}{E}\\
\parbox[t]{.43\linewidth}{$X$ \unsure{quasi-compact|quasicompact} (de degré $n$)} & \parbox[t]{.44\linewidth}{--- $E$ fini (\unsure{dé|de|d'é} ordre $n$)}\\[.5ex]
\parbox[t]{.43\linewidth}{\deleted{Décomposition de $X$ en}} & {}\\
\parbox[t]{.43\linewidth}{$X$ connexe [ou : irréductible], \inserted{ou $X\simeq\operatorname{Spec}\unsure{K|k'|K'}$, $K$ ext. finie sép.\unsure{]|)}}} & \parbox[t]{.44\linewidth}{--- $\pi$ transitif sur $E$}\\[.6ex]
\parbox[t]{.43\linewidth}{Décomposition de $X$ en somme de \unsure{compos.|composantes} \deleted{\unsure{disjointes|disj.}} connexes $\ne$} & \parbox[t]{.44\linewidth}{--- décomposition de $E$ en sommes d'orbites}\\[.6ex]
\parbox[t]{.43\linewidth}{$\varprojlim$ finies de schémas} & \parbox[t]{.44\linewidth}{--- $\varprojlim$ finies de $\pi$-ens}\\
\parbox[t]{.43\linewidth}{$\varinjlim$ \deleted{\unsure{fin|finies}} de schémas.} & \parbox[t]{.44\linewidth}{--- $\varinjlim$ de $\pi$-ens.}\\
\parbox[t]{.43\linewidth}{\unsure{mono|momo} $X\to Y$ \inserted{(resp. épi)}} & \parbox[t]{.44\linewidth}{--- $\unsure{E-F|E\to F}$ injectif (resp. surjectif)}\\
\parbox[t]{.43\linewidth}{[$X$ un torseur \inserted{à dr.} sous $G$ \deleted{\unsure{groupe f|groupe}} groupe fini ordinaire} & \parbox[t]{.44\linewidth}{$E$ un torseur à droite sous $G$; les opérations de $G$ commutant à l'action de $\pi$ (donc si on choisit un \unsure{$e\in E$|$e$ de $E$}, la structure est définie par un hom continu $\pi\to G$)]}
\end{array}\right.\)\endgroup}
\uline{Rappel sur la construction de} $\pi_0(X/k)$ lorsque \deleted{\unsure{pour des extensions|pour les extensions}} \deleted{\unsure{pour|génér.}} $X_{k'}$ ($k'$ clôture parfaite de $k$) est localement connexe, \deleted{\unsure{donc|dans}} \inserted{et} ses \deleted{\unsure{$X_{k'}$|$X_k$}} composantes connexes définissables sur des \inserted{\unsure{ss.|sép.|sé.}} ext. finies de $k$ (p. ex. $X$ qu. cpct, ou $X$ loc. de \uline{\unsure{pr. finie|p. finie|t. f.|t. fini}} sur $k$)

\mspage{2}
% archive sheet number: 9
\scanpage{11}
\deleted{\illegible{1}} Les pts \inserted{$x$} de $\underline{\pi_0}(X/k)$ correspondent aux \uline{composantes} \uline{connexes} de $X$, pour toute telle composante $X_{(x)}$, $k(x)$ est la plus grande sous-algèbre entière séparable de $\Gamma(X_{(x)},\underline{O}_{X_{(x)}})$ sur $k(x)$.

Que se passe-t-il pour $X=\operatorname{Spec}k'$ ? \deleted{\unsure{de|du}}

On a $k'=\varinjlim k_i$,\inserted{(\unsure{sous-ext.|sous ext.} galoisiennes finies de $k'$)} $\Rightarrow X=\varprojlim X_i$, $X_i=\operatorname{Spec} k_i$.

\[
\begin{aligned}
X'&=\varprojlim X'_i=\varprojlim_i(\pi_i)_{\unsure{k|k'}}=(\pi)_k\\
\pi_0(X')&=\pi\quad(\text{isom. \uline{topologique}})
\end{aligned}
\]
\deleted{Ce n'est plus un schéma étale sur $k$}

Schéma proétale sur $k$ [\deleted{\unsure{compact|complet}, réduit, à corps résiduels} \inserted{$=$ entier séparable $=$} \deleted{\unsure{à points ext.|à points isolés}} réduit à ext. résiduelles \deleted{\unsure{étal|étal.}} alg. sép.]

\deleted{\unsure{\S|S}} [Schéma loc. proétale sur $k$]

\[
\left\{\begin{aligned}
\mathrm{\unsure{ProEt|Proét}}(k)&\cong\pi\text{-esp cpt tdisc.}\cong\mathrm{Alg}_{\text{\unsure{indét|indet|ind-et}}}(k)\\
\mathrm{Locproet}(k)&\cong\pi\text{-esp loc. cpt disc.}
\end{aligned}\right.
\]
Si $X$ qu. cpct sur $k$ (pour simplifier) on construit une algèbre \unsure{proétale|indétale} \deleted{sur $X$} $k$, au choix

a)

\[
\begin{aligned}
A&=\text{limite inductive des sous-}\underline{\unsure{\mathrm{Alg}^{\prime\prime}\mathrm{bres}|\mathrm{Al}''\mathrm{bres}}}\ \text{étales de }\Gamma(X,\underline{O}_X)\\
&=\text{\deleted{limite} plus grande sous-Alg. ind-étale de }\Gamma(X,\underline{O}_X).\\
X&\longrightarrow\operatorname{Spec}(A)\overset{\text{\unsure{dfn|defn}}}{=}\underline{\pi}_0(X/k)
\end{aligned}
\]
b) Pour toute ext. \deleted{finie} \inserted{étale} $k_i\subset k'$ de $k$, on désigne par \deleted{\unsure{$k$|$A$}} $\unsure{A_i|A_{i'}}$ l'anneau des fonctions loc. constantes sur $X_{k_i}$, à valeurs dans $k_i$. Soit $\overline A=\varinjlim_i \unsure{{A_i}_{k'}|A_{i\,k'}}$, \inserted{fonctions loc. constantes sur $X_{k_i}$ \unsure{à val|à v.} dans $k'$} \deleted{\unsure{alors|ainsi} $\pi$ opère} \inserted{fonctions} loc. constantes sur $X_{k'}$ à valeurs dans $k'$, Alors \inserted{$\overline A$ est ind-étale et} $\pi$ opère dessus par transfert des structures, on pose $\overset{\deleted{\unsure{\text{alors}|\text{donc}}}}{A}=(\overline A)^\pi$.

\par\noindent\rule{.20\linewidth}{.4pt}\par
\[
\left\{\begin{array}{ll}
\text{a)} & \parbox[t]{.82\linewidth}{$\deleted{\unsure{X\to|X-}}\underline{\pi}_0(X/k)$ est proétale et $X\to\underline{\pi}_0(X/k)$ \deleted{est} surj. à fibres géom. connexes}\\[.7ex]
\text{b)} & \underline{\pi}_0(X/k)\ \text{fonctoriel en }X\\[.7ex]
\text{c)} & \underline{\pi}_0(X/k)\ \text{commute au produit}\\[.7ex]
\text{d)} & \underline{\pi}_0(X/k)\ \text{commute à l'extension de la base ....}
\end{array}\right.
\]
\uline{NB} $\underline{\pi}_0(X/k)$ peut se définir comme schém. loc proétale, \deleted{\unsure{si|dès}} \unsure{toute|tout} pt de \deleted{\unsure{ce|un}} $X$ a un voisinage ouvert et fermé $U$, tel que \deleted{\unsure{la|le}} la décomposition de $U_{k'}$ en somme \deleted{\unsure{de|s}} \deleted{\unsure{soit}} soit finie (p. ex. $U$ qu. cpct).

% archive sheet number: 10
\scanpage{12}
\uline{Schémas sur un corps.}

\deleted{\uline{Rappels} \textcircled{\unsure{1|2}} \unsure{Lorsque|Lorsque le} $k$ est parfait, \unsure{réduites.|réduite.}}

\margin{\textcircled{1} \uline{Séparabilité}}

% slanted
\uline{Rappel} Soit $k$ un corps. Conditions équivalentes

\[
\left\{\begin{array}{ll}
\text{a)}&\parbox[t]{.80\linewidth}{$k$ parfait i.e. $k$ de car. $0$, ou (si car $p\unsure{>|\ne}0$) $x\mapsto x^p$ est bijectif)}\\
\text{b)}&\text{un produit tensoriel d'algèbres réduites sur }k\text{ est}\\
\text{b')}&\Rightarrow\ \text{Un produit sur }k\text{ de schémas réduits est \inserted{\unsure{réduit|réduite}}}\\
\text{c)}&\parbox[t]{.80\linewidth}{Pour \unsure{tt|tte} extension finie $k'$ de $k$, \inserted{(\deleted{\unsure{quad|qu}} radicielle de degré $p$ si on veut)} \deleted{\unsure{on|un}} $k' \otimes_k k'$ est réduit}
\end{array}\right.
\]
Soit $A$ une algèbre sur un corps $k$, \inserted{$k_1$ une extension parfaite de $k$,} Propriétés équivalentes

\inserted{a) $\forall$ $k$ algèbre réduite $B$, $A\otimes_k B$ réduit}

b) $\forall$ extension $k'$ de $k$, $A\otimes_k k'$ réduit

c) $A\otimes_k k_1$ est réduit

d) $\forall$ ext. finie \deleted{\unsure{sép.|sép}} radicielle de $k$, $A\otimes_k k'$ est réduit.

\margin{faire passage \unsure{à la|à l'} limite pour prouver que $A\otimes_k k^{\unsure{\mathrm{pf}|\mathrm{perf}|p}}$ est réduit}

% slanted
\uline{Déf} On dit alors que $A$ est \uline{séparable} sur $k$ (ou \uline{géom.} \uline{réduit})

\uline{Cor} Soit $X$ schéma sur $k$. Propriétés équiv.

a) \deleted{\unsure{Pour tout|Pour tout un}} \unsure{tout|\deleted{tout}} \inserted{$k$-}schéma réduit $Y$, $X\times_k Y$ réduit

b) Pour tte \deleted{corps} extension $k'$ de $k$, $X_{k'}$ réduit

c) $X_{k_1}$ est réduit

\margin{\deleted{Cor}} d) (\unsure{Si|Soit} $X=\bigcup X_i$, les $X_i$ affines \deleted{\unsure{$\exists$|$A$}} d'anneaux $A_i$) les $A_i$ sont séparables $/k$

\uline{Déf} On dit alors que $X$ est sép. sur $k$, ou encore \deleted{\unsure{géo|géom.}} \inserted{géom.} réduit.

Ainsi, \deleted{$k$ est parfait \unsure{ssi|si}} tt schém. \inserted{géom.} réduit sur $k$ est \deleted{\unsure{géom.|géo}} réduit, la réciproque étant vraie si $k$ est parfait

\uline{NB} Si $k'/k$ est une extension de $k$, $\deleted{\unsure{X_f|X_{k'}|X_k}}X$ sép. sur $k$ $\Longleftrightarrow X_{k'}$ sép $/k'$

C'est une notion ``géométrique''. \hfill\uline{Stabilité par produit}

\margin{\textcircled{2} \deleted{\uline{\unsure{Intégrité|Intégralité}}} \uline{Irréductibilité}}

% slanted
\uline{Rappel} Soit $k$ \deleted{\unsure{sép clos|séparable}}, corps. Conditions équivalentes

a) $k$ sép. clos i.e. \deleted{\unsure{toute|tout}} ext. finie séparable de $k$ est $\unsure{\simeq|=} k$

b) Si $A$ et $B$ sont deux $k$-\deleted{\unsure{schémas irréductibles|schémas intègres}} \inserted{algèbres intègres}, $A\otimes B/\mathrm{\unsure{Nil|N.il}}$ est intègre

[b') Si $X,Y$ sont deux \inserted{$k$-}schémas irréd, alors $X\times_k Y$ est irréd]

c) Si \deleted{\unsure{$k'$|$k$}}, $k'$ est une extension finie de $k$, $k'\otimes_k k'\longrightarrow\deleted{\unsure{k|k'}}k'$ a un noyau nilpotent

\margin{Toute $\varprojlim$ filtrante d'espaces top. sobres irréd (resp. connexes) est irréd (resp. connexe)}

% slanted
Soit $X$ schém. sur corps, $k_1$ une extension sép. close de $k$ \uline{Conditions} équivalentes

a) $\forall Y$ sur $k$, $Y$ irréductible, $X\times_k Y$ est irréductible

b) $\forall$ extension $k'$ de $k$, $X_{k'}$ est irréductible

c) $X_{k_1}$ est irréductible

d) p. \unsure{toute|tte} ext. finie \deleted{\unsure{galois|gal.|galois.}} $k'$ de $k$, $X\otimes_k k'$ irréd.

\margin{faire passage à la limite ($\varprojlim$ d'irréd. par des morphismes dominants est irréd)}

% slanted
% archive sheet number: 11
\scanpage{13}
\uline{Déf} On dit alors que $X$ est géom. irréductible.

Ainsi \deleted{$k$ est sép. clos ssi} tout \inserted{$k$-}schéma irréductible est géom. irréductible, réciproque vraie ssi $k$ est sép. clos.

\uline{NB} Si $k'/k$ une extension, $\unsure{X/k|X_k/k}$ géom irréd $\Longleftrightarrow X_{k'}/k'$ géom irréd. \margin{$\big/$\,Stabilité par produit}

\margin{\textcircled{3}} \deleted{\uline{\unsure{Remarque|Réciproque}}}

\noindent\vrule width .55pt\hspace{.5em}\begin{minipage}{.92\linewidth}
\uline{Th} Soit $k$ un corps sép. clos. Si $X$ et $Y$ sont deux $k$-schémas connexes, $X\times_k Y$ est connexe.

\end{minipage}\par
\uline{Dém.} dans le cas où $X$ et $Y$ sont loc. de \unsure{t.|v.} f. sur $k$. (Cas général utilise technique de \unsure{descente|descents}\unsure{)|\}|?)}, plus précisément que $X\times Y\longrightarrow Y$ la top. de $Y$ est quotient de celle de $X\times Y$.

a) Si $X$ est connexe \inserted{non vide}, $X_{k'}$ est connexe \inserted{non vide}, pour tte extension $k'$ de $k$ [en effet, les fibres de $X_{k'}-X$ sont connexes \inserted{non vides}]

b) les fibres de $X\times Y\to Y$ étant connexes \inserted{non vides} et $Y$ connexe, $X\times Y$ est connexe

\uline{NB} Si $k$ est un corps tel que pour tt $X$ \inserted{connexe,} fini sur $k$, $X\unsure{\otimes|\times}_k \unsure{X|\alpha|x}$ est connexe alors $k$ est sép clos.

\uline{Prop} Soit $k$ un corps, $k_1$ une ext. \deleted{\unsure{alg|al}} sép. close, $X$ un $k$-schéma. Conditions équivalentes

a) $\forall$ $k$-schéma $Y$, $X\times_k Y$ est connexe

b) $\forall$ \unsure{ext.|lxt.} $k'$ de $k$, $X_{k'}$ est connexe

c) \deleted{$\unsure{\forall|\forall k}$} $X_{k_1}$ est connexe.

d) $\forall$ ext. finie \deleted{\unsure{gal.|gál.}} $k'$ de $k$, $X_{k'}$ est connexe.

\margin{passage à la limite, prenant $\unsure{k_1|k_0}=$ clôture sép \unsure{en $k$|en $k_1$|de $k$}, et en utilisant que $\operatorname{Spec} k_1 -\operatorname{Spec} k_{\unsure{i|0}}$ est \unsure{univ.|même} ouvert.}

% slanted
\uline{Déf.} On dit alors que $X$ est \uline{géom. connexe} sur $k$.

Ainsi \deleted{\unsure{connexe|connexe ssi}} \inserted{géom} connexe sur $k$ $\unsure{\Rightarrow|=}$ connexe, la réciproque vraie ssi $k$ est sép. clos.

\uline{NB} $X_k$ géom connexe $/k$ $\unsure{\Longleftrightarrow|\Longrightarrow} X_{k'}$ géom. connexe sur $k'$.

Stabilité par produit.

\margin{\textcircled{4} \uline{\unsure{poncticité|ponticité|ponctivité|intégrité}}}

% slanted
$k$ un corps. Conditions équivalentes

\deleted{a) Si $X$ et $Y$ deux $k$-algèbres intègres}

a) $k$ alg clos

b) Si $A,B$ sont deux $k$-algèbres intègres, $A\otimes_k B$ intègre

b') \deleted{\unsure{Si|Si les}} $X,Y$ intègres $\Rightarrow X\times_k Y$ \unsure{int.|int}

c) Pour tte ext. finie $k'$ de $k$, $k'\otimes_k k'$ intègre.

% archive sheet number: 12
\scanpage{14}
\uline{Prop} Soit $X$ un schéma sur corps $k$, $k_1$ une ext. alg. close de $k$. Cond. équiv.

a) $\deleted{\unsure{V|\forall}}\forall$ $k$-schéma intègre $Y$, $X\times_k Y$ intègre.

b) $\forall$ \unsure{ext.|lxt.} $k'$ de $k$, $X_{k'}$ intègre

c) \deleted{$\unsure{\forall|\forall k}$} $X_{k_1}$ intègre

\uline{Déf} On dit alors que $X$ est géom intègre ($\Longleftrightarrow$ géom irréd et géom réduit)

Ainsi géom intègre $\Rightarrow$ intègre, réciproque vraie ssi $k$ alg clos.

$\unsure{X_k/k|X/k}$ géom int. $\Longleftrightarrow X_{k'}/k'$ géom. intègre

Stabilité par \unsure{produits|produit}

\par\noindent\rule{.22\linewidth}{.4pt}\par
\uline{Prop} Soit $X$ sur corps $k$, $Z\xrightarrow{u}X$, $X$ connexe \inserted{\deleted{(irréductible)}}, $Z$ géom. connexe \inserted{non vide \deleted{(\unsure{resp.|resp} géom irréd)}}, (p. ex \inserted{non vide} \unsure{radiciel|rationnel}$/k$).

\noindent\vrule\enspace Alors $X$ est \deleted{\unsure{géom.|géomt}} connexe \deleted{(géom irréd)}.

Soit $k'$ \deleted{la clôture séparable de $X$.} \inserted{extension finie galoisienne}. Alors $\Pi=\operatorname{Gal}(k'/k)$ opère sur $X_{k'}=X\times_k\operatorname{Spec}(\unsure{k'|k'_i})$ et sur $Z_{k'}$, de façon à commuter à $u_{k'}: Z_{k'} - X_{k'}$.

\deleted{Les} \inserted{Soit $U$} parties ouvertes et fermées de $X_{k'}$, \deleted{\unsure{si il est|s'il est}} \inserted{prouvons qu'elle} \deleted{\unsure{est $\varnothing$, considérons|qu'on, considérons}} \inserted{est $\varnothing$ ou $X_{\unsure{k|k'}}$} prouvons qu'elle est stable sous l'action de $\Pi$. En effet,

\deleted{\unsure{Si|Soit} $g\in$ $U$ contient} Quitte à rempl. $U$ par $\unsure{\complement U|\int U}$, on peut supposer que $U$ contient $u_{k'}(Z_{k'})$, qui est stable par $\Pi$, donc l'intersection \inserted{\unsure{$\forall$|de|$\pi$}} des

\deleted{ouverts fermés qui contiennent $u_{\unsure{k_1|k'}}(Z_{k'})$ est \deleted{\unsure{bien|un}} stable sous $\Pi$,} \unsure{transformés,|transformés} de $U$ par les $g\in \unsure{\Pi|\mathrm{G}}$ est un ouvert stable \deleted{tel que} : \deleted{donc c'est l'image inverse d'une partie de $X$} non vide, provenant donc d'un ouvert \deleted{\unsure{stable de|et fermé}} de $X$, \unsure{etc|etc.}

\par\noindent\rule{.94\linewidth}{.4pt}\par
\uline{Ex} Composante neutre \inserted{$G^{\unsure{0|\circ}}$} d'un groupe algébrique \deleted{[elle est quasi-compacte} [Quand elle est ouverte, p. ex. $G$ loc. de \unsure{t. f.|t. f} sur $S$, alors on la munit de la structure induite]

$G_{\mathrm{red}}$ pour $k$ \uline{parfait} : un sous-groupe fermé de $G$

$G^{0}_{\mathrm{red}}$ id

[\unsure{Soit|Soient} $U$ \deleted{\unsure{et|deux}} \inserted{$\unsure{V|\varnothing}$} ouverts de $G$ ; alors $U.V^{-1}\supset G^0$ [\,$\Rightarrow$ et si $U,V$ \unsure{contient|contiennent} l'\unsure{élémt|élément} neutre, $UV\supset G^0$] [\unsure{Cor.|Car.} $G^0$ quasi-compact.]

% archive sheet number: 13
\scanpage{15}
Schémas (finis) étales sur un corps $k$ : la catégorie est équivalente \unsure{:|à} la catégorie des $\pi$-ensembles (finis), $\pi=\operatorname{Gal}(\bar k/k)$.

\uline{Application} Soit $X$ un schém. sur $k$, Supposons d'abord que $X_{\bar k}$ est loc. connexe [p. ex. loc. noeth, ou ens des composantes irréductibles loc. fini ; \unsure{tout est OK|tout est ok} si $X$ est loc. de t.f. sur $k$]. Alors $\unsure{I|\bar I}=\pi_0(X_{\bar k})$ est un ens. discret, et on a un \deleted{\unsure{f|$\varphi$}} morphisme de $\bar k$-schémas

\diagram{\begin{tikzcd}[ampersand replacement=\&]
X_{\bar k}\arrow[d]\\
I_{\bar k}
\end{tikzcd}}
(il est universel pour les morphismes de $X$ dans des schémas constants). D'ailleurs $\pi$ opère sur $\deleted{\unsure{I_{\bar k}|X_{\bar k}}} I$, donc $I_{\bar k}$ provient d'un \deleted{\unsure{$I$ sur|truc}} schém. étale $T$ sur $k$ [si on prouve que $\pi$ opère continûment -- ce sera le cas en tous cas si $X$ est quasi-compact].

\margin{i.e \deleted{\unsure{tte|toute}} composante connexe de $X_{\bar k}$ est \unsure{déjà|déja} définie sur une \inserted{\unsure{ss-|s.-|sép.}}ext. finie de $k$. \uline{C'est vrai} \uline{si} $X$ \uline{qu. cpct}, \deleted{\unsure{mais|aussi}} \deleted{\illegible{2}}}

% slanted
Alors la théorie de la descente nous prouve que $X\simeq X_{\bar k}/\pi$, $T=I_{\bar k}/\pi$, (dans catégorie des schémas) donc on trouve

\diagram{\begin{tikzcd}[ampersand replacement=\&]
X\arrow[d,"f"]\\
T
\end{tikzcd}}
$T$ étant étale sur $k$, \deleted{\unsure{et|etc}} On \unsure{a|â|à} que $f$ est surjectif et à fibres géom connexes -- et ceci suffit à caractériser $X$. On peut aussi décrire comme le morphisme universel de $X$ dans un schém. étale sur $k$.

\margin{$T$ est appelé le schém. des composantes connexes géom. de $X$. Un pt à val. dans un $S'/k$ ``est'' une partie ouverte et fermée de $X_{\unsure{S'|S}}$, qui \unsure{est|\deleted{est}}}

\diagram{\begin{tikzcd}[ampersand replacement=\&,column sep=large]
X\arrow[d] \& X'\arrow[d,dash]\\
T\arrow[d] \& T'\arrow[d,dash]\\
S=\unsure{\operatorname{Spec}|\operatorname{Sp}} k\arrow[r,dash] \& S'
\end{tikzcd}}
\uline{Généralisation} On suppose $X$ quasi-compact, ou somme de schémas quasi \unsure{cpcts|cpts} : On trouve un $T$ \inserted{affine} \unsure{pro étale|proétale}, ayant les \unsure{\^{m}s|\~{m}s|mêmes} propriétés\unsure{$\scriptscriptstyle\wedge$|,|$\scriptscriptstyle\times$} qu'avant. Le foncteur qu'il représente est un peu plus subtil à décrire....

% archive sheet number: 14
